% ---------------------------------------------------------------
% 1. Introdução
% ---------------------------------------------------------------
\chapter{Introdução}\label{cap:intro}

\section{Contexto e motivação}\label{sec:contexto}

O presente projeto tem como objetivo o desenvolvimento de um sistema pericial de apoio à decisão, em condução autónoma, no qual o conhecimento do domínio é representado, explicitamente, através de factos, regras e prioridades e avaliado por dois motores de inferência: \emph{Drools} e \emph{Prolog}. Uma vez que este relatório está inserido e foi desenvolvido para a unidade curricular de ENGCIA, apenas será abordado o motor de inferência no âmbito do \emph{Drools}.

O sistema a apresentar destina-se a fabricantes de automóveis que tenham como objetivo investir em soluções de condução autónoma. No entanto, não basta que o veículo seja capaz de tomar a ação correta, é também necessário demonstrar o porquê da decisão ter sido tomada e o porquê de ter rejeitado outras alternativas, seja para validar o comportamento em cenários de teste, seja para responder às exigências da circulação do dia-a-dia, de homologação e de segurança [5, 6].

Neste sentido, a aplicação recebe continuamente os dados dos sensores simulados, atualiza a memória de trabalho e avalia as regras sobre o estado do veículo e do ambiente. Uma vez que várias regras podem ser aplicadas no mesmo ciclo e indicar ações diferentes, o sistema utiliza dois canais independentes para determinar a decisão final: o da velocidade, em que prevalecem as situações de emergência e de paragem obrigatória e são respeitadas as restrições de velocidade; e o da via, em que a decisão é selecionada de acordo com os critérios de prioridade definidos. Cada decisão guarda não só a regra, a justificação e a evidência que lhe deu origem, como também as propostas que foram descartadas e o motivo de não terem sido selecionadas. Como o foco do presente trabalho não é representar um veículo autónomo completo, mas sim estudar o ciclo de decisão, é possível rastreá-lo através da perceção, da inferência, da tomada de decisão, da atuação e, por fim, de uma nova observação.

\section{Problemática e questões orientadoras}\label{sec:problematica}

O problema abordado consiste em transformar as observações dos sensores, como a velocidade, a distância ao veículo da frente, a presença de peões, a sinalização e as condições meteorológicas, em factos para serem utilizados pelo motor de inferência. Alguns destes factos podem manter-se válidos durante várias leituras ou sofrer alterações entre leituras consecutivas como, por exemplo, o limite de velocidade ou a distância ao veículo da frente, respetivamente. Com base nesta informação, o sistema avalia, em cada ciclo, as condições de circulação e aplica as regras definidas para determinar as ações a executar, tendo em conta os critérios de prioridade estabelecidos. Neste sentido, o processo de inferência e tomada de decisão pode ser representado pela Figura 1.

\begin{figure}[htbp]
\centering
\begin{tikzpicture}[
    node distance=6mm and 7mm,
    caixa/.style={draw=corEscura, line width=0.8pt, rounded corners=3pt, fill=corFundo,
                  align=center, font=\small, minimum height=13mm, text width=24mm},
    seta/.style={-{Stealth[length=2.4mm]}, line width=0.8pt, draw=corEscura}
]
\node[caixa] (lei) {Leitura dos\\sensores};
\node[caixa, right=of lei] (fac) {Eventos e factos\\na memória\\de trabalho};
\node[caixa, right=of fac] (mot) {Motor de\\inferência:\\avaliação\\das regras};
\node[caixa, right=of mot] (pro) {Propostas de\\decisão};
\node[caixa, below=10mm of pro] (sel) {Seleção da\\decisão segundo\\as prioridades};
\node[caixa, left=of sel] (aca) {Ação sobre o\\veículo};
\node[caixa, left=of aca] (nov) {Novo estado\\do veículo e\\do ambiente};
\draw[seta] (lei) -- (fac);
\draw[seta] (fac) -- (mot);
\draw[seta] (mot) -- (pro);
\draw[seta] (pro) -- (sel);
\draw[seta] (sel) -- (aca);
\draw[seta] (aca) -- (nov);
\draw[seta] (nov.west) -| node[pos=0.25, below, font=\scriptsize, text=corCinza]{nova leitura} (lei.south);
\end{tikzpicture}
\caption{Processo de inferência e tomada de decisão do sistema pericial}
\label{fig:inferencia}
\end{figure}

A partir do problema apresentado, foi definida a seguinte questão para orientar o desenvolvimento
do projeto:

\emph{De que forma pode um motor de inferência utilizar, continuamente,} dados provenientes de diferentes sensores para a tomada de decisão, \emph{tendo em conta a evolução temporal da informação?}

\noindent Desta questão decorrem as seguintes questões orientadoras:

\begin{itemize}
  \item De que forma podemos utilizar os dados provenientes de diferentes sensores para manter
        uma representação atualizada da situação do veículo e do ambiente?
  \item Até que ponto a análise de várias leituras ao longo do tempo pode ajudar a distinguir uma
        situação pontual de uma situação de perigo persistente?
  \item De que forma é que o sistema pode selecionar uma decisão quando várias regras são aplicáveis em simultâneo, apesar de apresentarem propostas de ações diferentes?
  \item Como pode o sistema explicar o processo que o levou a uma determinada conclusão ou ação,
        através da identificação dos factos utilizados e das respetivas regras (\ing{How})?
  \item É possível explicar o motivo pelo qual uma determinada decisão foi tomada e o porquê de não terem sido escolhidas outras alternativas (\emph{Why} e \emph{Why Not})?
\end{itemize}

\section{Objetivos}\label{sec:objetivos}

Para concretizar o objetivo geral do projeto, foram definidos os seguintes objetivos específicos:

\begin{itemize}
  \item Representar o conhecimento adquirido através de factos, eventos e regras de produção em
        \ing{Drools}, para permitir interpretar as diferentes situações de circulação;
  \item Desenvolver mecanismos de inferência que permitam analisar os dados provenientes dos
        sensores simulados e a respetiva evolução temporal, para identificar situações de risco e
        determinar as ações de condução adequadas;
  \item Definir critérios de prioridade e resolução de conflitos, de forma a selecionar as
        decisões de velocidade e de via quando várias regras são aplicáveis em simultâneo;
  \item Integrar diferentes mecanismos de raciocínio, nomeadamente o processamento de eventos
        complexos, o encadeamento regressivo e a lógica difusa, para analisar situações temporais,
        verificar a segurança das manobras e adaptar a velocidade às condições meteorológicas;
  \item Assegurar a explicabilidade das decisões, para permitir identificar os factos e as regras
        do processo de inferência, bem como justificar as decisões tomadas e as alternativas
        rejeitadas;
  \item Avaliar o comportamento do sistema através dos cenários de simulação para verificar a coerência das decisões e o funcionamento dos mecanismos de explicação.
\end{itemize}
